\SourceLeaf{63}
by $\left(1-\dfrac{r}{\nu}\right)$, and since multiplication by $\dfrac{r}{\nu}$ can be carried out by
differentiating the preceding term with respect to $n$, multiplying it again by $n$ and dividing by $\nu\,\mathrm{d}n$,
we obtain the following relation for the initial term of $B^{(r)}$
\[
= A^{(r)}\text{ term. init.} - \frac{n\,{}^{*}\!\mathrm{d}A^{(r)}\text{ term. init.}}{\nu\,\mathrm{d}n}.
\]
Exactly the same relation holds for the general terms. For the coefficient
$A^{(r)}$, the term at index $h^{\mathrm{tus}}$ is
\[
= \frac{1}{2^{r+2h-1}}\,{}^{r+2h}\mathfrak{B}^{(r+h)} . {}^{\nu}\mathfrak{B}^{(r+2h)}\,n^{r+2h}.
\]
Thus, on changing $\nu$ to $\nu-1$, the term of $B^{(r)}$ at index $h^{\mathrm{tus}}$ becomes
\[
= \frac{1}{2^{r+2h-1}}\,{}^{r+2h}\mathfrak{B}^{(r+h)} . {}^{\nu-1}\mathfrak{B}^{(r+2h)}\,n^{r+2h},
\]
which follows directly from the preceding expression by introducing just the factor
\[
\frac{\nu-(r+2h)}{\nu} = 1 - \frac{r+2h}{\nu}
\]
into it. Now, multiplying the term of
$A^{(r)}$ at index $h^{\mathrm{tum}}$ by $\dfrac{r+2h}{\nu}$ is equivalent to differentiating it with respect to
$n$, dividing by $\nu\,\mathrm{d}n$ and multiplying again by $n$. We therefore have
\[
B^{(r)}\text{ term. }h = A^{(r)}\text{ term. }h - \frac{n\,{}^{*}\!\mathrm{d}A^{(r)}\text{ term. }h}{\nu\,\mathrm{d}n},
\]
and since this relation holds for every term without exception, it gives the equation
\[
B^{(r)} = A^{(r)} - \frac{n\,{}^{*}\!\mathrm{d}A^{(r)}}{\nu\,\mathrm{d}n}.
\]
\par\noindent\hspace*{\parindent}$\varepsilon$) By this recurrence, $(1+n\cos\varphi)^{\nu-1}$ is thus reduced to $(1+n\cos\varphi)^{\nu}$
in such a way that the coefficients of the former can be derived very easily by differentiation
from the coefficients that occur in the latter expression. Since the coefficients of $(1+n\cos\varphi)^{-1}$
are available in finite form, we can proceed step by step to negative exponents of greater magnitude
and thus express negative integral powers in general by finite coefficients.

The infinite series found earlier for these coefficients are nevertheless useful.
They apply whenever $\nu$ is a fraction. In that case, in \textsc{Euler}'s judgment, which carries great
weight in these investigations, there is no prospect of obtaining finite formulas
